% Figure from MATH 551 notes, section 10. Compile with the notes preamble.
\begin{tikzpicture}[scale=1.0]
  \def\Epath{plot[smooth cycle, tension=0.8] coordinates {(0,0.2) (1.2,-0.5) (2.6,-0.2) (3.0,1.0) (2.5,2.3) (1.1,2.6) (0.1,1.7)}}
  \def\Tpath{plot[smooth cycle, tension=0.8] coordinates {(1.9,0.6) (3.0,0.1) (4.4,0.5) (4.7,1.4) (3.9,2.1) (2.6,1.9) (1.8,1.4)}}
  \fill[figB!8] \Epath;
  \begin{scope}
    \clip \Tpath;
    \fill[figR!18] (-1,-1) rectangle (6,4);
    \fill[figB!40] \Epath;
  \end{scope}
  \draw[figB, thick] \Epath;
  \draw[black!75, thick] \Tpath;
  \node[font=\small, figB] at (0.8,1.2) {$E$};
  \node[font=\small, black!80] at (4.35,2.2) {$T$};
  \node[font=\scriptsize, figB!60!black] at (2.45,1.05) {$T \cap E$};
  \node[font=\scriptsize, figR!80!black] at (3.75,1.05) {$T \cap E^c$};
  \node[font=\scriptsize, black!60] at (4.3,-0.35) {$E^c$};
\end{tikzpicture}
